%%============================================================
\subsection{Data Mining}
%%============================================================

% Giới thiệu ngắn: Tensor decomposition trong data mining —
% mục tiêu không chỉ nén mà còn trích xuất thành phần có ý nghĩa
% (interpretable components) từ dữ liệu đa chiều.
% Liên kết ngược với Chương 2 (cơ sở lý thuyết CP, Tucker).

% TODO: Viết đoạn mở đầu (~3–5 câu)

%%------------------------------------------------------------
\subsubsection{Tính duy nhất CP và khả năng diễn giải}
\label{sec:dm-uniqueness}
%%------------------------------------------------------------

% Nội dung cần viết:
% - Nhắc lại Định lý Kruskal (Định lý~\ref{thm:kruskal}):
%   điều kiện k_A + k_B + k_C >= 2R + 2
% - Hệ quả: CP decomposition cho kết quả duy nhất (essentially unique)
%   mà KHÔNG cần ràng buộc trực giao — khác biệt cốt lõi so với SVD/PCA
% - So sánh trực tiếp với SVD (Bảng~\ref{tab:cp-vs-svd}):
%   SVD cần rotation → mất ý nghĩa vật lý
%   CP không cần rotation → mỗi component có ý nghĩa trực tiếp
% - Ý nghĩa thực tiễn: trong data mining, ta cần components
%   có ý nghĩa vật lý (spatial pattern, temporal pattern, subject grouping)
%   — không chỉ cần variance explained
% - Ví dụ minh họa: tensor EEG (Subjects × Channels × Time)
%   → mỗi CP component cho 3 vectors diễn giải được

% TODO: Viết nội dung (~0.5–0.75 trang)

%%------------------------------------------------------------
\subsubsection{Thuật toán CP-ALS và CP-OPT}
\label{sec:dm-algorithms}
%%------------------------------------------------------------

% === CP-ALS (Alternating Least Squares) ===
% Nội dung cần viết:
% - Ý tưởng: cố định N-1 factor matrices, tối ưu factor matrix còn lại
% - Bước cập nhật (sử dụng ký hiệu từ Chương 2):
%   A^(n) ← MTTKRP(X, factors, n) · (Hadamard of Gram matrices)^{-1}
% - MTTKRP: phép toán cốt lõi, tính hiệu quả bằng einsum
% - Chuẩn hóa cột + hấp thụ norm vào weights λ
% - Tiêu chí hội tụ: relative fit change < ε
% - Ưu điểm: đơn giản, mỗi bước có closed-form solution
% - Nhược điểm: hội tụ chậm (ALS swamp), chỉ tuyến tính

% TODO: Viết phần CP-ALS (~0.5 trang)
% TODO: (Tùy chọn) Thêm pseudocode Algorithm 1

% === CP-OPT (Gradient-Based, L-BFGS-B) ===
% Nội dung cần viết:
% - Ý tưởng: vector hóa toàn bộ factor matrices → 1 vector x ∈ R^p
% - Hàm mục tiêu: f(x) = ||X - X̂||²_F (tính hiệu quả qua Gram matrices)
% - Gradient: ∂f/∂A^(n) = -2·MTTKRP + 2·A^(n)·V_n
% - Sử dụng L-BFGS-B (scipy): quasi-Newton, xấp xỉ Hessian
% - So sánh với ALS:
%   + Hội tụ nhanh hơn gần nghiệm (superlinear vs linear)
%   + Tốn kém hơn mỗi iteration (full gradient)
%   + Robustness hơn với ALS swamp

% TODO: Viết phần CP-OPT (~0.5 trang)

% === So sánh CP-ALS vs CP-OPT ===
% TODO: Thêm bảng so sánh (iterations, fit, thời gian, FMS)
% \begin{table}[H]
% \centering
% \begin{tabular}{lccc}
% \hline
% \textbf{Thuật toán} & \textbf{Iterations} & \textbf{Fit} & \textbf{FMS} \\
% \hline
% CP-ALS & & & \\
% CP-OPT & & & \\
% \hline
% \end{tabular}
% \caption{So sánh CP-ALS và CP-OPT.}
% \label{tab:als-vs-opt}
% \end{table}

%%------------------------------------------------------------
\subsubsection{PARAFAC2 và CMTF}
\label{sec:dm-parafac2-cmtf}
%%------------------------------------------------------------

% === PARAFAC2 ===
% Nội dung cần viết:
% - Động lực: dữ liệu thực thường có kích thước không đều
%   (ví dụ: số trial khác nhau giữa các subject trong EEG)
% - Mô hình: X_k ≈ P_k · diag(s_k) · V^T
%   với P_k thay đổi theo slice k, V chung
% - Ràng buộc cross-product: P_k^T P_k = Φ (hằng số qua k)
%   → đảm bảo tính duy nhất mà vẫn linh hoạt
% - Thuật toán PARAFAC2-ALS:
%   project → stack → ALS update → Procrustes rotation
% - Ưu điểm so với CP thông thường:
%   xử lý irregular data mà không cần padding/truncation

% TODO: Viết phần PARAFAC2 (~0.5 trang)

% === CMTF (Coupled Matrix-Tensor Factorization) ===
% Nội dung cần viết:
% - Động lực: tận dụng dữ liệu bổ sung (side information)
%   Ví dụ: tensor EEG (Subjects × Channels × Time) kết hợp
%   với ma trận metadata (Subjects × Demographics)
% - Mô hình: phân rã tensor X và ma trận Y đồng thời,
%   chia sẻ factor matrix trên mode chung (Subjects)
% - Hàm mục tiêu: min ||X - X̂||² + α||Y - Ŷ||²
% - Lợi ích: cải thiện chất lượng phân rã nhờ regularization
%   từ dữ liệu bổ sung; kết nối các nguồn dữ liệu khác nhau

% TODO: Viết phần CMTF (~0.25–0.5 trang)

%%------------------------------------------------------------
\subsubsection{Ứng dụng}
\label{sec:dm-applications}
%%------------------------------------------------------------

% === Ứng dụng 1: Phân tích EEG (Motor Imagery) ===
% Nội dung cần viết:
% - Bộ dữ liệu: PhysioNet EEG Motor Movement/Imagery
%   (109 subjects, 64 channels, 160 Hz)
% - Tensor: Subjects × Channels × Time (sau tiền xử lý)
% - Tiền xử lý: bandpass 8–30 Hz, epoch extraction, baseline correction
% - Kết quả CP decomposition:
%   + Chọn rank (elbow plot + CORCONDIA)
%   + Diễn giải 3 modes: subject loadings, channel topography, temporal signature
%   + So sánh CP vs SVD: CP tách biệt spatial–temporal, SVD trộn lẫn
% - PARAFAC2 trên EEG: xử lý khi số trial khác nhau giữa subjects

% TODO: Viết phần EEG (~0.5–0.75 trang)
% TODO: Thêm figures từ notebooks (elbow plot, factor heatmaps, topography maps)

% === Ứng dụng 2: Metabolomics / Gut Microbiome ===
% Nội dung cần viết:
% - Bối cảnh: dữ liệu metabolomics thường có dạng
%   Samples × Metabolites × Time (hoặc Conditions)
% - Tại sao tensor: tương tác nhiều chiều (sample–metabolite–condition)
%   không thể phân tích bằng PCA 2D thông thường
% - CP uniqueness đặc biệt hữu ích: mỗi component tương ứng
%   với một biological pathway hoặc metabolic profile
% - Ví dụ: PARAFAC2 cho chromatography data (elution profiles
%   có chiều dài khác nhau giữa các mẫu)
% - Liên hệ tutorial NeurIPS 2025: metabolomics là một trong
%   ba ứng dụng chính được nhấn mạnh

% TODO: Viết phần Metabolomics (~0.25–0.5 trang)
% Có thể tham khảo: Bro (1997), Acar et al. (2007)